\documentclass[10pt, a4paper]{article}

\usepackage[utf8]{inputenc}
\usepackage[spanish]{babel}
\usepackage{geometry}
\geometry{top=1.8cm,bottom=1.8cm,left=2cm,right=2cm}
\usepackage{booktabs}
\usepackage{array}
\usepackage{enumitem}
\usepackage{amsmath}
\usepackage{tikz}
\usetikzlibrary{shapes.geometric, arrows.meta, positioning}
\usepackage{hyperref}

\title{\textbf{Análisis de Caso: Fallo en Sistema Distribuido Real\\AWS us-east-1 (2011)}}

\author{
Universidad Técnica Estatal de Quevedo\\
Facultad de Ciencias de la Computación y Diseño Digital\\
Asignatura: Aplicaciones Distribuidas
}

\date{Semana 3 --- Unidad 1}

\begin{document}

\maketitle
\vspace{-1cm}

\section{Descripción del Incidente}

El 21 de abril de 2011, Amazon Web Services (AWS) sufrió una interrupción crítica en la región \textit{us-east-1}, afectando servicios como Amazon EC2 y Amazon RDS durante aproximadamente 96 horas. El incidente comenzó debido a un cambio de red mal ejecutado que desvió tráfico hacia una red de alta latencia en lugar de la red de baja latencia destinada para EC2.

Como consecuencia, miles de volúmenes EBS (\textit{Elastic Block Store}) detectaron pérdida de réplicas y comenzaron simultáneamente un proceso automático de re-replicación (\textit{re-mirroring}). La sobrecarga causada saturó el almacenamiento disponible, generando un efecto de retroalimentación positiva conocido como \textbf{``re-mirroring storm''}.

\subsection*{Impacto Cuantificado}

\begin{itemize}[noitemsep]
    \item Servicios afectados: EC2, RDS, Elastic Beanstalk, Heroku, Reddit, Quora y Foursquare.
    \item Aproximadamente el 1.87\% de instancias EC2 quedaron inaccesibles.
    \item El 0.07\% de volúmenes EBS sufrió pérdida permanente de datos.
    \item Duración de la interrupción severa: 96 horas.
    \item Pérdidas económicas estimadas: decenas de millones de dólares.
\end{itemize}

\section{Análisis Técnico}

\subsection{Clasificación del Fallo}

\begin{center}
\small
\begin{tabular}{p{3cm} p{9cm}}
\toprule
\textbf{Tipo de fallo} & \textbf{Presencia en el incidente} \\
\midrule
Fallo de omisión & Las instancias EC2 dejaron de responder a solicitudes. \\
Fallo de temporización & La alta latencia provocó tiempos de respuesta críticos. \\
Fallo de crash & Volúmenes EBS quedaron en estado ``stuck''. \\
Fallo bizantino & No se presentaron respuestas incorrectas. \\
\bottomrule
\end{tabular}
\end{center}

\subsection{Violación del Teorema CAP}

AWS EBS estaba diseñado bajo un modelo \textbf{CP} (Consistencia + Tolerancia a Particiones). Ante la pérdida de réplicas, el sistema decidió preservar la consistencia bloqueando el acceso a los datos hasta restaurar todas las réplicas. Como consecuencia, la disponibilidad del sistema cayó prácticamente a cero durante la recuperación.

\subsection{Transparencias ANSA Comprometidas}

\begin{itemize}[noitemsep]
    \item Transparencia de fallos.
    \item Transparencia de replicación.
    \item Transparencia de disponibilidad.
    \item Transparencia de ubicación.
\end{itemize}

\section{Diagrama de Ishikawa --- Causa Raíz}

\begin{center}

\begin{tikzpicture}[
scale=0.85,
transform shape,
every node/.style={font=\small},
box/.style={
draw,
rounded corners,
align=center,
fill=gray!10,
minimum width=3cm,
minimum height=0.9cm
}
]

% Línea principal
\draw[very thick,->] (0,0) -- (10,0);

% Cabeza
\node[box, minimum width=3.2cm, minimum height=1cm]
at (10.8,0)
{\textbf{Caída AWS}\\us-east-1};

% Espinas superiores
\draw[thick] (2,0) -- (0.8,1.5);
\draw[thick] (5,0) -- (3.8,1.5);
\draw[thick] (8,0) -- (6.8,1.5);

% Espinas inferiores
\draw[thick] (2,0) -- (0.8,-1.5);
\draw[thick] (5,0) -- (3.8,-1.5);
\draw[thick] (8,0) -- (6.8,-1.5);

% Causas superiores
\node[box] at (-0.3,1.9)
{\textbf{Error humano}\\Cambio incorrecto\\de red};

\node[box] at (2.8,1.9)
{\textbf{Red}\\Tráfico enviado\\a red lenta};

\node[box] at (5.8,1.9)
{\textbf{Latencia}\\Altos tiempos\\de respuesta};

% Causas inferiores
\node[box] at (-0.3,-1.9)
{\textbf{Replicación}\\Re-espejado\\masivo};

\node[box] at (2.8,-1.9)
{\textbf{Sobrecarga}\\Saturación del\\almacenamiento};

\node[box] at (5.8,-1.9)
{\textbf{Control}\\Ausencia de\\\textit{rate limiting}};

\end{tikzpicture}

\end{center}

\textbf{Causa raíz identificada:} ausencia de mecanismos de
\textit{rate limiting} en los procesos automáticos de recuperación,
combinada con un error humano en la configuración de red.

\section{Lecciones Aprendidas}

\begin{itemize}[noitemsep]
    \item Implementar \textit{rate limiting} y \textit{Circuit Breaker}.
    \item Diseñar arquitecturas Multi-AZ.
    \item Realizar pruebas de resiliencia (\textit{Chaos Engineering}).
    \item Definir explícitamente si el sistema será CP o AP.
    \item Aplicar procedimientos estrictos de \textit{change management}.
\end{itemize}

\section{Aplicación al Proyecto Final}

Para el PFC se propone:

\begin{enumerate}[noitemsep]
    \item Implementar \textit{Circuit Breaker} entre microservicios.
    \item Incorporar \textit{health checks} con \textit{timeouts}.
    \item Ejecutar pruebas de tolerancia a fallos con Docker.
    \item Documentar decisiones arquitectónicas relacionadas con CAP.
\end{enumerate}

\section{Conclusión}

El fallo de AWS us-east-1 demostró que los sistemas distribuidos no son tolerantes a fallos únicamente por tener múltiples nodos. La resiliencia depende del diseño correcto de los mecanismos de recuperación, monitoreo y control de errores. Además, evidenció cómo un pequeño error operativo puede amplificarse y provocar una caída masiva cuando no existen límites adecuados en los procesos automáticos del sistema.

\vspace{0.3cm}

\section*{Referencias}

\small

[1] Amazon Web Services, ``Summary of the Amazon EC2 and Amazon RDS Service Disruption in the US East Region'', 2011. Disponible en: \url{https://aws.amazon.com/message/65648/}

[2] E. Brewer, ``Towards robust distributed systems'', PODC, 2000.

[3] S. Gilbert and N. Lynch, ``Brewer's conjecture and the feasibility of consistent, available, partition-tolerant web services'', ACM SIGACT News, 2002.

\end{document}
```